\hspace{3mm}

\centerline{\bf \Huge 8 класс}

\begin{flushright}
        \scriptsize{}
\end{flushright}

\problem{\textbf{1}}
Докажите, что нет ни одного четырёхзначного палиндрома, делящегося на 27.

\problem{\textbf{2}}
В ряд стоят четыре человека, среди которых есть Веня (который всегда говорит правду), Федя (который всегда лжёт) и преподаватели Расул Владимирович с Анджуром Аркадьевичем (которые могут сказать как правду, так и солгать).

--- Самый левый сказал: <<Рядом со мной стоит Веня>>.

--- Второй слева сказал: <<Рядом со мной стоит Федя>>.

--- Третий слева сказал: <<Рядом со мной стоят преподаватели>>.

--- Самый правый сказал: <<Рядом со мной стоит Веня>>.

Определите, кто и где может стоять.

\problem{\textbf{3}}
Анджур Аркадьевич открыл кафе <<ЭлМШ>>, повесив на входе табличку с надписью <<только для Шиншилл и Лососей>>. Лера считала процентное соотношение Шиншилл и после некоторого притока посетителей их стало больше 98\% от общего числа. Правда ли, что в какой-то момент их было ровно 98\%, если первым зашёл Лосось?

\problem{\textbf{4}}
Трапеция $A B C D$ обладает следующими свойствами: $A B \| C D, A B=3 C D$ и $C D=D A$. Найдите углы трапеции, зная, что $\angle C D A=120^{\circ}$.

\problem{\textbf{5}}
У Элианы есть $19$ монет, одна из которых фальшивая и весит меньше настоящих, а все настоящие весят одинаково. Имеются также двое двухчашечных весов. Нелли испортила весы, и теперь они ломаются, если одна из чаш перевесит другую (однако можно успеть понять, какая именно чаша перевесила). Сможет ли она за три взвешивания определить фальшивую монету (возможно, поломав в итоге весы)?